%
% crosshost_connection_lifecycle.tex
%
% W12 (lux-833p), Invariant 1 of the cross-host transport design (DES-090):
% "no content before verification."  Models the connection lifecycle
% (accepted -> TLS handshaking -> TLS verified, awaiting ConnectMessage ->
% identified, HubId cross-checked -> receiving content) and proves that no
% SceneMessage is processed against a connection that has not passed BOTH
% gates.
%
% Models the code in:
%   src/punt_lux/display/cross_host_listener.py    -- CrossHostListener
%       (Gate 1: the pending-handshake set; a socket is not in the reader
%        set until do_handshake() has returned "ready")
%   src/punt_lux/display/cross_host_verification.py -- CrossHostVerification
%       (Gate 2: the SAN/hub_id cross-check)
%   src/punt_lux/display/hub_reconciliation.py     -- handle_connect
%       (Gate 2: reject-or-register-identity)
%   src/punt_lux/display/identity_guard.py         -- reject_scene_unless_hub
%       (content requires an identified kind="hub" fd)
% and the invariant stated in
%   docs/architecture/system.tex "Cross-Host Transport and Trust",
%   subsection "Invariants", item 1, and subsection "Coexistence with the
%   Local Fast Path" ("Two distinct gates, not one").
%
% Type-check:  fuzz docs/crosshost_connection_lifecycle.tex
% Model-check (goal predicates in the Verification section):
%   probcli docs/crosshost_connection_lifecycle.tex -model_check \
%       -p DEFAULT_SETSIZE 3
%
% Formatting follows the fuzz house rules: \quad~ for continuation indent,
% schema lines under 80 columns, predicates broken at \land/\lor/\implies.
%

\documentclass[a4paper,10pt,fleqn]{article}
\usepackage[margin=1in]{geometry}
\usepackage{fuzz}
\usepackage[colorlinks=true,linkcolor=blue,citecolor=blue,urlcolor=blue]{hyperref}

\begin{document}

\title{Cross-Host Connection Lifecycle: a Z Specification}
\author{W12 formal model of DES-090 Invariant 1 (no content before
verification)}
\date{September 2026}
\maketitle

% ============================================================
\section{Introduction}
% ============================================================

The same-host Hub-to-Display leg is an \field{AF\_UNIX} socket whose
\texttt{0700} permission \emph{is} the trust model: anything that can open
the path is already a process the same OS user controls.  Cross-host
removes that free lunch --- anyone who can route a packet to the
Display's port can attempt to speak the Hub protocol --- so a
newly-connected peer must clear two gates before its content is honored,
and the accept path must never let a \msg{SceneMessage} slip past either
one.

A connection's life runs through four observable conditions:

\begin{description}
  \item[accepted.] The TCP connection exists; the socket sits in a
    pending-handshake set, deliberately \emph{not} in the reader set, so
    \texttt{poll\_clients} never reads a frame from it
    (\texttt{CrossHostListener.accept\_pending}).
  \item[TLS verified (Gate 1).] \texttt{do\_handshake()} has returned
    successfully --- mutual certificate verification complete.  Only now
    is the fd promoted into the ordinary reader set, where it may exchange
    \msg{ConnectMessage} like any other connection
    (\texttt{CrossHostListener.pump\_ready}).
  \item[identified (Gate 2).] A \msg{ConnectMessage} arrived and its
    self-reported \field{hub\_id} hostname matched the peer certificate's
    SAN (\texttt{CrossHostVerification}); \texttt{handle\_connect}
    recorded the fd as \texttt{kind="hub"}.  A mismatch instead closes the
    connection, never a degraded-trust fallback.
  \item[receiving content.] Only an identified \texttt{kind="hub"} fd has
    its \msg{SceneMessage} processed (\texttt{reject\_scene\_unless\_hub}).
\end{description}

The hazard the design closes is an accept path that hands a socket to the
ordinary reader set --- making its frames deliverable --- before its TLS
handshake, or before its \msg{ConnectMessage}, has completed, letting a
\msg{SceneMessage} land against an unverified or not-yet-identified
connection.  The handshake itself is driven non-blocking, one step per
frame under a bounded deadline, precisely so a stalled peer cannot freeze
the render thread; but the correctness property this model concerns is
the \emph{ordering}, not the timing: content is honored only after both
gates.

Three properties must hold across every interleaving:

\begin{description}
  \item[I1 --- no identify before the handshake (Gate 1).] No connection
    is identified until its TLS handshake has completed.  Because content
    requires identification (I2), this is what keeps a \msg{SceneMessage}
    from an un-handshaken peer out.
  \item[I2 --- no content before identification (Gate 2).] No
    \msg{SceneMessage} is processed against a connection that has not
    been identified (SAN-cross-checked \texttt{kind="hub"}).
  \item[I3 --- deadlock freedom.] Some operation is always enabled.
\end{description}

Together I1 and I2 are exactly Invariant 1: content only after \emph{both}
transport verification and \field{hub\_id} cross-check.

% ============================================================
\section{Basic Types}
% ============================================================

\begin{zed}
  [CONN]
\end{zed}

\texttt{CONN} is the set of cross-host connections (file descriptors) the
Display can hold at once; two suffice to exhibit one connection racing
through the gates while another is at a different stage, and three give
ProB room to interleave a stalled, a verified, and an identified
connection.  The carrier is bounded by ProB's \texttt{DEFAULT\_SETSIZE}.

% ============================================================
\section{State}
% ============================================================

The state is flat: one schema of five subsets of \texttt{CONN}, each
recording a condition a connection has reached.  A connection may hold
several at once (an identified one is also handshaken and readable); the
subset relations that say so are I1 and I2, and are deliberately
\emph{not} written into the schema.

\begin{schema}{Conn}
  known : \power CONN \\
  readable : \power CONN \\
  handshook : \power CONN \\
  identified : \power CONN \\
  served : \power CONN
\where
  readable \subseteq known \\
  handshook \subseteq known \\
  identified \subseteq known \\
  served \subseteq known
\end{schema}

\texttt{known} is every connection currently open (accepted, not yet
dropped).  \texttt{readable} is the ordinary reader set ---
\texttt{SocketListener}'s client dict, the membership \texttt{poll\_clients}
tests before it ever delivers a frame to a handler; a connection outside
it can have no message of any kind processed.  \texttt{handshook} is
Gate 1 complete.  \texttt{identified} is Gate 2 complete.  \texttt{served}
records every connection against which at least one \msg{SceneMessage} has
been processed --- the observable this model exists to constrain.  The
four containments are structural hygiene: dropping a connection removes it
from every set, so nothing outlives \texttt{known}.  The \emph{ordering}
containments --- \texttt{identified \subseteq handshook} (I1) and
\texttt{served \subseteq identified} (I2) --- are the safety properties,
checked as reachability of their negations (Section~\ref{sec:verify}), not
enforced as guards here.

% ============================================================
\section{Initialization}
% ============================================================

\begin{schema}{Init}
  Conn'
\where
  known' = \emptyset \\
  readable' = \emptyset \\
  handshook' = \emptyset \\
  identified' = \emptyset \\
  served' = \emptyset
\end{schema}

A single initial state: no connection exists --- the cold start before any
peer has connected to the cross-host listener.

% ============================================================
\section{Operations}
% ============================================================

\subsection{Accept --- into the pending set, not the reader set}

A peer completes the TCP handshake.  \texttt{c?} becomes known but is
\emph{not} made readable: it sits in the pending-handshake set, kept
separate from the reader set exactly so no frame can be read from it until
its TLS handshake is verified complete.  This is the ordering the whole
design turns on --- \texttt{accept\_pending} wraps the raw socket
non-blocking with \field{do\_handshake\_on\_connect=False} and holds it
apart from \texttt{\_clients}/\texttt{\_readers}.

\begin{schema}{Accept}
  \Delta Conn \\
  c? : CONN
\where
  c? \notin known \\
  known' = known \cup \{ c? \} \\
  readable' = readable \\
  handshook' = handshook \\
  identified' = identified \\
  served' = served
\end{schema}

\subsection{HandshakeComplete --- Gate 1, promote to the reader set}

\texttt{do\_handshake()} returns successfully for a pending connection.
Only now is \texttt{c?} both recorded as handshaken and promoted into the
reader set --- \texttt{pump\_ready} returning the socket and the caller
moving it out of pending into \texttt{\_clients}.  The two happen together:
there is no state in which a connection is readable but not handshaken.

\begin{schema}{HandshakeComplete}
  \Delta Conn \\
  c? : CONN
\where
  c? \in known \\
  c? \notin handshook \\
  handshook' = handshook \cup \{ c? \} \\
  readable' = readable \cup \{ c? \} \\
  known' = known \\
  identified' = identified \\
  served' = served
\end{schema}

\subsection{Identify --- Gate 2, SAN match}

A \msg{ConnectMessage} arrives on a readable connection and its
\field{hub\_id} hostname matches the peer certificate's SAN.
\texttt{handle\_connect} records the fd as identified.  The guard is
\texttt{c? \in readable}: the message could only be delivered because the
fd is in the reader set.  It does \emph{not} re-check \texttt{handshook}
--- Gate 1's ordering is what guarantees a readable fd is already
handshaken, and removing that ordering (Section~\ref{sec:fidelity},
variant A) is exactly how a not-yet-verified connection reaches identify.

\begin{schema}{Identify}
  \Delta Conn \\
  c? : CONN
\where
  c? \in readable \\
  c? \notin identified \\
  identified' = identified \cup \{ c? \} \\
  known' = known \\
  readable' = readable \\
  handshook' = handshook \\
  served' = served
\end{schema}

\subsection{SanReject --- Gate 2, SAN mismatch, fail-closed}

The other outcome of Gate 2: the \msg{ConnectMessage}'s \field{hub\_id}
does not match the certificate SAN, so \texttt{handle\_connect} closes the
connection (\texttt{remove\_client}) rather than identifying it --- never a
degraded-trust fallback (Invariant 3).  \texttt{c?} leaves every set.

\begin{schema}{SanReject}
  \Delta Conn \\
  c? : CONN
\where
  c? \in readable \\
  c? \notin identified \\
  known' = known \setminus \{ c? \} \\
  readable' = readable \setminus \{ c? \} \\
  handshook' = handshook \setminus \{ c? \} \\
  identified' = identified \\
  served' = served \setminus \{ c? \}
\end{schema}

\subsection{ProcessScene --- content, only past both gates}

A \msg{SceneMessage} is processed.  The guard is the conjunction of both
gates: \texttt{c? \in readable} (the frame was delivered) \emph{and}
\texttt{c? \in identified} (Gate 2 passed) --- the two facts
\texttt{reject\_scene\_unless\_hub} checks before a scene is installed.
This guard \emph{is} I2; weakening it to \texttt{readable} alone
(Section~\ref{sec:fidelity}, variant B) is the W1 gap where an
unidentified fd's scene is installed.

\begin{schema}{ProcessScene}
  \Delta Conn \\
  c? : CONN
\where
  c? \in readable \\
  c? \in identified \\
  served' = served \cup \{ c? \} \\
  known' = known \\
  readable' = readable \\
  handshook' = handshook \\
  identified' = identified
\end{schema}

\subsection{Drop --- an ordinary disconnect}

The connection's socket closes on its own.  \texttt{c?} leaves every set
--- the existing per-fd reaping, applied identically to a cross-host fd.

\begin{schema}{Drop}
  \Delta Conn \\
  c? : CONN
\where
  c? \in known \\
  known' = known \setminus \{ c? \} \\
  readable' = readable \setminus \{ c? \} \\
  handshook' = handshook \setminus \{ c? \} \\
  identified' = identified \setminus \{ c? \} \\
  served' = served \setminus \{ c? \}
\end{schema}

% ============================================================
\section{Invariants}
\label{sec:inv}
% ============================================================

\paragraph{Invariant 1 --- no identify before the handshake (Gate 1).}
\[
  identified \subseteq handshook.
\]
\texttt{Identify} requires \texttt{c? \in readable}, and \texttt{readable}
is granted only by \texttt{HandshakeComplete}, which grants
\texttt{handshook} in the same step --- so \texttt{readable \subseteq
handshook} holds throughout, and hence so does \texttt{identified
\subseteq handshook}.  \texttt{Accept} adds to \texttt{known} only, never
to \texttt{readable}; that is the ordering.  The check
(Section~\ref{sec:verify}) asks whether the negation --- an identified
connection that is not handshaken --- is reachable.

\paragraph{Invariant 2 --- no content before identification (Gate 2).}
\[
  served \subseteq identified.
\]
\texttt{ProcessScene} requires \texttt{c? \in identified} before recording
it in \texttt{served}, and no other operation grows \texttt{served}.  With
I1, this gives \texttt{served \subseteq identified \subseteq handshook}:
content only past both gates --- Invariant 1 of the design in full.

\paragraph{Invariant 3 --- deadlock freedom.}
\texttt{Accept} is enabled for any \texttt{c? \notin known} whenever one
exists; in the sole state where it has no fresh argument
(\texttt{known = CONN}), \texttt{Drop} is enabled.  So some operation is
always enabled.

% ============================================================
\section{Verification}
\label{sec:verify}
% ============================================================

Type-checking is a \texttt{fuzz} pass:
\begin{verbatim}
  fuzz docs/crosshost_connection_lifecycle.tex
\end{verbatim}

I1 and I2 are state properties, checked by asking ProB whether their
negations are reachable.  A goal reported \emph{not found} means the
invariant holds on every reachable state.  Both return \emph{goal not
found} against \texttt{DEFAULT\_SETSIZE 3}:
\begin{verbatim}
  # Invariant 1: an identified-but-not-handshaken connection
  #              (goal NOT found for the correct spec)
  probcli docs/crosshost_connection_lifecycle.tex -model_check -nodead \
    -goal "card({c | c : CONN & c : identified & not(c : handshook)}) > 0" \
    -p DEFAULT_SETSIZE 3

  # Invariant 2: a served-but-not-identified connection
  #              (goal NOT found for the correct spec)
  probcli docs/crosshost_connection_lifecycle.tex -model_check -nodead \
    -goal "card({c | c : CONN & c : served & not(c : identified)}) > 0" \
    -p DEFAULT_SETSIZE 3
\end{verbatim}

I3 is the full-state-space deadlock check.  Every state has an enabled
operation (Section~\ref{sec:inv}), so any deadlock reported would be a
model defect:
\begin{verbatim}
  # Invariant 3 / deadlock freedom (no deadlock, all states visited)
  probcli docs/crosshost_connection_lifecycle.tex -model_check \
    -p DEFAULT_SETSIZE 3
\end{verbatim}

The witnesses must each be \emph{found}, or the invariants would hold only
because the mechanism they guard is unreachable.  Both return \emph{goal
found} for the correct spec:
\begin{verbatim}
  # W1: content is actually processed (goal FOUND)
  probcli docs/crosshost_connection_lifecycle.tex -model_check -nodead \
    -goal "card(served) > 0" \
    -p DEFAULT_SETSIZE 3

  # W2: the between-gates window is reachable -- handshaken (Gate 1 done)
  #     but not yet identified (Gate 2 pending) (goal FOUND)
  probcli docs/crosshost_connection_lifecycle.tex -model_check -nodead \
    -goal "card({c | c : CONN & c : handshook & not(c : identified)}) > 0" \
    -p DEFAULT_SETSIZE 3
\end{verbatim}

% ============================================================
\section{Fidelity: reproducing the shipped defect}
\label{sec:fidelity}
% ============================================================

A model that cannot reproduce the known defect proves nothing.  Two buggy
variants, one per gate, each remove exactly one ordering guarantee and
each surfaces its own invariant violation.

\paragraph{Variant A --- handshake-before-accept removed (Gate 1).}
Kept as
\texttt{docs/crosshost\_connection\_lifecycle\_no\_handshake\_gate\_buggy.tex},
this is the concrete race the design names: an accept path that adds a
socket to the reader set before its TLS handshake has completed.
\texttt{Accept} becomes
\begin{verbatim}
  % correct (this spec): accepted != readable
  readable' = readable
  % buggy (variant A): freshly accepted, immediately readable
  readable' = readable \cup {c?}
\end{verbatim}
Now a \msg{ConnectMessage} can be delivered --- and \texttt{Identify} can
fire --- on a connection whose TLS handshake never completed.  Trace:
\texttt{Accept}(c) makes \texttt{c} readable while
\texttt{c \notin handshook}; \texttt{Identify}(c) fires (guard
\texttt{c \in readable} holds); \texttt{c} is now identified but not
handshaken.  I1's negation is reachable: its goal returns \emph{found}
against variant A and \emph{not found} against the correct spec.  This is
an unauthenticated peer being honored as an identified Hub --- the mTLS
handshake and its SAN check bypassed entirely.

\paragraph{Variant B --- identify gate removed (Gate 2, the W1 gap).}
Kept as
\texttt{docs/crosshost\_connection\_lifecycle\_no\_identify\_gate\_buggy.tex},
this is the shipped same-host gap W1 closes: a content handler that
accepts a \msg{SceneMessage} from any readable fd, not only an identified
one.  \texttt{ProcessScene}'s guard becomes
\begin{verbatim}
  % correct (this spec):
  c? : readable & c? : identified
  % buggy (variant B):
  c? : readable
\end{verbatim}
Trace: \texttt{Accept}(c); \texttt{HandshakeComplete}(c) --- \texttt{c}
readable and handshaken but not yet identified; \texttt{ProcessScene}(c)
fires (the weakened guard does not require \texttt{identified}).  I2's
negation is reachable: its goal returns \emph{found} against variant B and
\emph{not found} against the correct spec.  This is a scene installed
under a connection with no verified \texttt{HubId} to attribute it to ---
harmless behind the \texttt{0700} \field{AF\_UNIX} permission, a hole once
the store is keyed by \texttt{HubId}.

Each variant isolates one gate: variant A leaves I2's guard intact (so
I2 still holds there --- the failure is purely I1), and variant B leaves
Gate 1 intact (so I1 still holds there --- the failure is purely I2).
Together they show the model detects a breach of \emph{either} gate, which
is why Invariant 1 of the design requires \emph{both}.

% ============================================================
\section{Test Partitions}
\label{sec:partitions}
% ============================================================

Derived from the operations and invariants above, for audit against the
test suite (see \texttt{docs/crosshost\_transport\_coverage.md}):

\begin{description}
  \item[R1 --- an accepted, un-handshaken connection is not readable.] A
    freshly accepted socket sits in the pending set; no frame is read from
    it.  The Gate 1 ordering.
  \item[R2 --- Gate 1 promotes on handshake success.] A completed
    handshake, and only then, moves the fd into the reader set.
  \item[R3 --- the between-gates window is real.] A handshaken,
    not-yet-identified connection is readable but content-refused (witness
    W2).
  \item[R4 --- Gate 2 identifies on SAN match.] A \msg{ConnectMessage}
    whose \field{hub\_id} matches the certificate SAN records the fd as
    \texttt{kind="hub"}.
  \item[R5 --- Gate 2 rejects, fail-closed, on SAN mismatch.] A mismatch
    closes the connection; it never becomes identified and never serves
    content.
  \item[R6 --- content only past both gates.] A \msg{SceneMessage} is
    processed only for an identified connection (I2), which is
    handshaken (I1) --- so served $\subseteq$ handshaken, model-checked
    exhaustively.
  \item[R7 --- an un-handshaken connection never serves content.] The
    Gate 1 regression against variant A.
  \item[R8 --- an unidentified connection never serves content.] The
    Gate 2 regression against variant B (the W1 gap).
  \item[R9 --- a drop at any stage removes the connection cleanly.] No
    residual membership in any set survives a disconnect.
\end{description}

\end{document}
